% ===== §1 Introduction: Generativity Under Power =====
\section{Introduction: Generativity Under Power}
\label{sec:intro}

\subsection*{Prelude: Two Bodies Falling}

\noindent
Two bodies fall toward one another across a dark field, and do not meet. Each is held to its own path by something it did not choose and cannot simply set down---a momentum gathered before the other was in view, a curvature laid into it by everything that made it what it is. They draw near; the nearness is real; and at the closest point of approach they pass, each bending the other's path a little, neither arriving. This is not a failure of the falling. It is the shape that falling takes between two bodies that are genuinely two.

We are taught to read such a passing as a wound---the near-miss, the word that lands wrong, the tenderness offered in a language the other cannot read, the silence where an answer was wanted. And sometimes it is a wound. But before it is anything to be grieved or repaired, the passing is the trace of something the falling could not have done without: that the two remained two, that neither was absorbed into the other's path, that there was, between them, a difference strong enough to bend a trajectory and not be dissolved by it. A love that admitted no such passing would be a love in which one body had captured the other entirely---and a captured body no longer falls. The Japanese have a word for this passing-without-meeting: \zh{すれ違い}, \emph{surechigai}. This paper is, in one of its several aspects, an extended inquiry into why it happens, and why its happening is not, in the first instance, a sign that something has gone wrong.

This paper is written in the knowledge that the field is not empty. Around the two bodies, and through them, run other forces: the pull of families who hold what the two still need; the scripts a society writes for how a life together ought to go; the slow sedimentation of two material histories into one shared and contested present. These forces are not external decorations upon a private bond. They enter it. They bend it. And the oldest hope of lovers---that if the love were only deep enough, only equal enough, only good enough, the forces would fall away and leave the two alone in a clearing of their own---is, this paper will argue, a hope that mistakes the nature of the field. The forces do not fall away. The question is not how to escape them, but how, beneath them and through them, two people might still create something that is theirs; how what they create might survive the forces; and how, in time, it might come to be free.

\goldrule

\subsection{The three questions}
\label{sub:intro-three}

This paper asks how love fares under power. It is the seventeenth in a series that has developed, across its earlier installments, an account of intimate relation as a \emph{generative-relational} process: a coupling of two subjects out of which value is not extracted but \emph{generated}, and through which, when the coupling is good, that value returns to its creators rather than being siphoned to a centre.\footnote{For the foundational diagnosis on which the series rests, see \citet{wan2024grb}. For the generative-relational framework and its formal criterion---the positive \emph{holonomy} of a good cycle, distinguished from the zero or negative holonomy of an extractive one---see \pinkref{Paper~IX} \citep{paperix}; for the political economy of relational wealth drawn on here, \pinkref{Papers~XIV--XVI}; and for the theory of relational trust that the present paper's account of trust extends, \pinkref{Paper~XIII} \citep{paperxiii}.} The series has given this generativity a formal criterion, an ontological ground---the \emph{relational subject}, constituted in and by its relations rather than prior to them---and a poetics: the good symbolic cycle as one that refuses settlement and invites its participants to walk a path and become witnesses. What the series has not yet confronted, directly and as its central problem, is \emph{power}: the plain fact that intimate relation does not occur in a clearing but in a field saturated with asymmetries, and that these asymmetries can deform, capture, freeze, and on occasion shatter the very generativity the series has laboured to describe.

The argument is organised around three questions, each answered by a different body of theory, and each having the same object---\emph{love}, in the precise sense the series gives that word: the dynamics of the real between two subjects, the part of their bond that circles the unsymbolisable and cannot be counted, traded, or directly managed.

The first question is \emb{how love creates a world}---how, out of the coupling of two subjects, anything genuinely new is generated at all. This is the question the framework of the generative-relational being (GRB) was built to answer, and the present paper takes its affirmative findings largely as established. What the paper adds, on this first front, is the \emph{etiology of its failure}: a systematic account of how the world-creating coupling breaks down---how estrangement, the passing-without-meeting of the Prelude, actually comes about. We will find that it comes about in many ways, and that they fall into a small number of layers.

The second question is \emb{how love is not destroyed by power}. Here the paper turns to the theory of international relations (IR)---not, as will be argued at length, to its realist core, which assumes that the fundamental currency of an anarchic field is \emph{relative power}, an assumption the generative-relational account must reject; but to a reconstructed IR whose true object is the more general one of \emph{how opaque subjects, unable to enter one another directly, build trust and order under uncertainty and without a central authority}. So defined, IR turns out to describe the intimate field with uncanny precision---more purely, in one respect, than it describes the international one. But it describes only the field's \emph{means}: the symbolic order of power through which a relation defends the conditions of its own flourishing. It cannot reach the \emph{end} those means serve.

The third question is \emb{how love becomes free within history}. Power structures are not given; they are the historical sediment of particular material conditions, and they are therefore, in principle, historically mutable. To see this---and to see it without sliding into a crude determinism that would reduce love itself to an effect of its economic base---the paper turns, in a dedicated section, to a \emph{dialectical} historical materialism: one that holds the determining force of material conditions together with the relative autonomy, and the reactive force, of the symbolic, ethical, and affective life those conditions shape but do not exhaust.

Three questions, then, with a single object and an ascending logic: \emb{creation, preservation, liberation}. Love must first be able to generate; what is generated must then be able to survive the forces around it; and survival is not yet enough---what survives must be able, in time and against the inertia of inherited structure, to become free. The three theories answer at three different distances, but they answer about the same thing; and none of them, it must be said at once, can touch that thing directly. This is the paper's central methodological commitment, and it requires its own statement.

\subsection{Two levels: the end and the means}
\label{sub:intro-twolevels}

The intimate relation, as this paper treats it, is a two-level system.

At the level of the \emb{end} stands the dynamics of love itself---a process of the \emph{real}, in the Lacanian register the series has used throughout. It circles a vacancy (\emph{das Ding}) that cannot be fully symbolised; it is not a quantity and cannot be measured; it is not a good and cannot be exchanged. It is the blooming, not the trellis.

At the level of the \emb{means} stands the symbolic order of power---the asymmetries, the defences, the alliances and balances and boundary-works through which a relation secures its place in a field of forces. This level \emph{is} calculable; it \emph{can} be theorised by IR and by the political economy of relational wealth the series has developed elsewhere; it is the proper domain of strategy.

The means serve the end. But---and this is the source of nearly every difficulty the paper will trace---the means have an \emph{intrinsic tendency to damage the end they serve}. A defence mounted to protect the love can, in the mounting, import into the bond the very logic of relative advantage the love cannot survive. The trellis can be built so tight that it strangles the bloom.

\begin{claim}[The political economy of conditions]
The tools of power and of wealth---all of which operate in the symbolic---can protect the \emph{conditions} under which love occurs, but they cannot produce or directly secure the love itself. For the dynamics of the real, touched directly by a symbolic instrument, do not survive the touch: managed, they decohere; counted, they are no longer what they were. One can build a fence to keep out what would trample the flower, but one cannot fence a flower into bloom. What the political economy of the intimate can offer is therefore, strictly, a \emph{political economy of conditions}---never a political economy of love.
\end{claim}

This is why \emph{trust} stands in the title beside power and estrangement. Trust is the one thing in the system that crosses both levels. It can be cultivated, in part, by symbolic means---by costly and credible signals, by the slow projection of a real fidelity into observable form---and it is also rooted in the real, in a coupling that cannot be faked. Trust is the interface: the bridge by which the political economy of conditions reaches toward, without ever seizing, the dynamics of love it exists to protect. It is also, as we will see, the variable that decides whether an asymmetry of power becomes generative or frozen---and so it runs through every question the paper asks.

\subsection{Two claims}
\label{sub:intro-claims}

Two claims organise what follows.

The \emb{first} concerns estrangement. The \emph{surechigai} of the Prelude---the structural passing-without-meeting---is not, in the first instance, a malfunction of the bond but a necessary moment in the bond's own self-maintenance. A relation that never passed, that achieved and held a perfect concordance, would be a relation at thermodynamic equilibrium: which is to say, dead. This was the thesis of the paper's earliest conception, and it survives here, but demoted from the whole to a special case---one outcome, at the structural-necessity end, of a far wider etiology of estrangement that the paper now lays out in full.

The \emb{second}, and the paper's centre of gravity, concerns power. It can be stated as a double negation enclosing a single positive. Power \emph{cannot eliminate itself}: there is no mechanism by which a sufficiency of love or equality or goodwill causes power to dissolve, and the belief that there is---call it the optimistic naturalism of intimacy---is false. Power also \emph{cannot be eliminated entirely}: it is the by-product of difference, and difference cannot be removed without removing, along with it, the very creativity that makes the relation generative---so that the radical utopian dream of a relation cleansed of power is not merely unattainable but self-defeating. Between these two negations lies the only viable course:

\begin{claim}[Generative balance]
To generate, alongside the good value a relation creates, a continual counter-force that keeps power from solidifying. A generative relation does not oppose power. It opposes \emph{the freezing of power}.
\end{claim}

\noindent This, the paper will argue at its theoretical summit (\xref{sec:balance}), is not an external ethical demand imposed upon generativity from without, but the political appearance of an ontological law that generativity already obeys---the same law by which desire, generating its symbols, withholds itself from total symbolisation in order to remain alive.

\subsection{Scope, and the road ahead}
\label{sub:intro-scope}

A word on what the paper is and is not. It is not a manual of relationship technique; the practices it will eventually describe are reconstructed as the \emph{maintenance of conditions}, never as the management of love. It is not a defence of any particular political arrangement; where its argument ascends, at one point, toward a general criterion for evaluating political bodies---whether they generate, alongside their value, the counter-force that keeps their own power from freezing---it marks that ascent as the seed of a separate work and does not pursue it here. And it is not a reduction of love to power, or of the bond to a calculus of forces; the entire architecture of the two levels exists precisely to hold open the space no calculus reaches.

The road runs as follows. A symptomatology (\xref{sec:symptom}) sets out the observable forms in which power damages generativity. The etiology proper (\xref{sec:encoding}--\xref{sec:ir}) traces estrangement through four lenses---the mismatch of trust-grammars (\xref{sec:encoding}), the unattributability of prediction error and the active withdrawal it triggers (\xref{sec:prediction}), the arrest of geometric-phase accumulation (\xref{sec:phase}), and the order-problem of opaque subjects that requires the reconstruction of IR (\xref{sec:ir}). The paper then turns to the means and their limits: the political economy of conditions and why it cannot reach the real (\xref{sec:conditions}); the inventory of classical IR practices and their teleological filtering (\xref{sec:practices}); and the two summits---the negative one, on generative balance and the suppression of power's freezing (\xref{sec:balance}), and the positive one, on the asymmetry that generates rather than dominates (\xref{sec:asymmetry}). A dedicated section (\xref{sec:histmat}) gives the dialectical-materialist history of intimate power. A cosmological movement (\xref{sec:dissipation}--\xref{sec:rose}) situates estrangement within the thermodynamics of generative systems, the grammar of return, and the dormancy of the rose. A praxis of the two axes (\xref{sec:praxis}) and a polyphonic conclusion with its envoi (\xref{sec:conclusions}) close the paper. Throughout, the discipline announced in the series holds: the frameworks are not to be synthesised into a master theory but held in their plurality and their dialogue, each speaking to the edge of what it can say, and no farther.
